\documentclass[aspectratio=169]{beamer}

\usetheme{Madrid}
\usecolortheme{default}

\usepackage{amsmath}
\usepackage{booktabs}
\usepackage{tikz}
\usepackage{graphicx}
\usepackage{hyperref}

\title{Modern Foundation Models}
\subtitle{Lecture 2: Transformer Architecture Details}
\author{Dror Lederman}
\institute{Holon Institute of Technology}
\date{\today}

\begin{document}

%------------------------------------------------
\begin{frame}
\titlepage
\end{frame}

%------------------------------------------------
\begin{frame}{Lecture Roadmap}
\begin{enumerate}
    \item Recap: Self-Attention and Transformer Blocks
    \item Position Embeddings
    \item Layer Normalization
    \item Attention Approximation
    \item Transformer-Based Model Families
    \item BERT and Encoder-Only Models
\end{enumerate}
\end{frame}

%------------------------------------------------
\section{Recap}

\begin{frame}{Recap: Scaled Dot-Product Attention}

The core attention operation is:

\[
\mathrm{Attention}(Q,K,V)
=
\mathrm{softmax}
\left(
\frac{QK^{T}}{\sqrt{d_k}}
\right)V
\]

\vspace{0.4cm}

\begin{itemize}
    \item $Q$ = queries
    \item $K$ = keys
    \item $V$ = values
    \item $d_k$ = key/query dimension
\end{itemize}

\end{frame}

%------------------------------------------------
\begin{frame}{Why Attention Matters}

Self-attention allows each token to directly attend to other tokens.

\vspace{0.4cm}

Example:

\begin{quote}
The law will never be perfect, but its application should be just.
\end{quote}

\vspace{0.3cm}

The token ``its'' may attend to:
\begin{itemize}
    \item ``law''
    \item ``application''
\end{itemize}

\vspace{0.3cm}

Attention can help model long-range dependencies.

\end{frame}

%------------------------------------------------
\begin{frame}{Multi-Head Attention}

Instead of using one attention mechanism, Transformers use multiple heads.

\[
\mathrm{MultiHead}(Q,K,V)
=
\mathrm{Concat}(\mathrm{head}_1,\ldots,\mathrm{head}_h)W^O
\]

\[
\mathrm{head}_i =
\mathrm{Attention}(QW_i^Q,KW_i^K,VW_i^V)
\]

\vspace{0.4cm}

Different heads may capture different linguistic or semantic relations.

\end{frame}

%------------------------------------------------
\section{Position Embeddings}

\begin{frame}{Why Do We Need Position Information?}

Self-attention is permutation-invariant.

\vspace{0.3cm}

Without position information:

\[
\text{``the dog bit the man''}
\]

and

\[
\text{``the man bit the dog''}
\]

may become too similar.

\vspace{0.4cm}

Transformers therefore need a way to represent token order.

\end{frame}

%------------------------------------------------
\begin{frame}{Learned Position Embeddings}

A simple approach:

\[
x_i = e_i + p_i
\]

where:

\begin{itemize}
    \item $e_i$ is the token embedding
    \item $p_i$ is a learned position embedding
\end{itemize}

\vspace{0.4cm}

Limitation:

\begin{itemize}
    \item Learned position embeddings may not extrapolate well to longer sequences.
\end{itemize}

\end{frame}

%------------------------------------------------
\begin{frame}{Sinusoidal Position Embeddings}

The original Transformer used fixed sinusoidal encodings:

\[
PE_{m,2i}
=
\sin
\left(
\frac{m}{10000^{2i/d_{\text{model}}}}
\right)
\]

\[
PE_{m,2i+1}
=
\cos
\left(
\frac{m}{10000^{2i/d_{\text{model}}}}
\right)
\]

\vspace{0.3cm}

where:
\begin{itemize}
    \item $m$ is the position
    \item $i$ is the embedding dimension index
\end{itemize}

\end{frame}

%------------------------------------------------
\begin{frame}{Why Sinusoids?}

Using trigonometric identities:

\[
\cos(a-b)
=
\cos(a)\cos(b)
+
\sin(a)\sin(b)
\]

\vspace{0.4cm}

This means that the similarity between position embeddings can encode relative distance:

\[
\langle PE_m, PE_n \rangle = f(m-n)
\]

\vspace{0.4cm}

So the model can reason about relative positions.

\end{frame}

%------------------------------------------------
\begin{frame}{From Absolute to Relative Position}

For attention, what often matters is not the absolute position:

\[
m
\]

but the relative distance:

\[
m-n
\]

\vspace{0.4cm}

Modern models often modify the attention mechanism directly to encode relative positions.

\end{frame}

%------------------------------------------------
\begin{frame}{Position Bias in Attention}

A relative position bias can be added to the attention score:

\[
\mathrm{softmax}
\left(
\frac{\langle q_m,k_n\rangle}{\sqrt{d_k}}
+
\mathrm{bias}(m,n)
\right)
\]

\vspace{0.4cm}

Examples:

\begin{itemize}
    \item T5: learned relative position bias
    \item ALiBi: deterministic linear bias
\end{itemize}

\end{frame}

%------------------------------------------------
\begin{frame}{ALiBi: Attention with Linear Biases}

ALiBi uses a simple distance-based bias:

\[
\mathrm{bias}(m,n) = \mu (n-m)
\]

\vspace{0.4cm}

Main idea:

\begin{itemize}
    \item Encourage attention to nearby tokens.
    \item Allow extrapolation to longer sequences.
    \item Avoid explicit position embeddings.
\end{itemize}

\end{frame}

%------------------------------------------------
\begin{frame}{RoPE: Rotary Position Embeddings}

RoPE rotates query and key vectors according to their positions.

For two dimensions:

\[
R_{\theta,m}
=
\begin{pmatrix}
\cos(m\theta) & -\sin(m\theta) \\
\sin(m\theta) & \cos(m\theta)
\end{pmatrix}
\]

\vspace{0.4cm}

Apply rotation to queries and keys:

\[
q_m \rightarrow q_m R_{\theta,m}
\]

\[
k_n \rightarrow k_n R_{\theta,n}
\]

\end{frame}

%------------------------------------------------
\begin{frame}{Why RoPE Works}

After rotation, the attention score depends on relative distance:

\[
q_m k_n^T
=
x_m W_q R_{\theta,n-m} W_k^T x_n^T
\]

\vspace{0.4cm}

Main advantage:

\begin{itemize}
    \item Relative position is naturally encoded.
    \item Works well in modern decoder-only LLMs.
    \item Common in Llama-style architectures.
\end{itemize}

\end{frame}

%------------------------------------------------
\section{Layer Normalization}

\begin{frame}{Layer Normalization}

Layer Normalization normalizes activations within a layer.

\[
\mathrm{LN}(x)
=
\gamma \hat{x} + \beta
\]

where:

\[
\hat{x}
=
\frac{x-\mu}{\sqrt{\sigma^2+\epsilon}}
\]

\[
\mu = \frac{1}{d}\sum_{i=1}^{d}x_i
\]

\[
\sigma^2 = \frac{1}{d}\sum_{i=1}^{d}(x_i-\mu)^2
\]

\end{frame}

%------------------------------------------------
\begin{frame}{Why Layer Normalization?}

Layer Normalization helps with:

\begin{itemize}
    \item Training stability
    \item Faster convergence
    \item Better gradient flow
    \item Training very deep networks
\end{itemize}

\vspace{0.4cm}

In Transformers, LayerNorm is used together with residual connections.

\end{frame}

%------------------------------------------------
\begin{frame}{Post-Norm vs Pre-Norm}

\textbf{Post-Norm:}

\[
x_{l+1}
=
\mathrm{LayerNorm}
\left(
x_l + \mathrm{SubLayer}(x_l)
\right)
\]

\vspace{0.5cm}

\textbf{Pre-Norm:}

\[
x_{l+1}
=
x_l + \mathrm{SubLayer}
\left(
\mathrm{LayerNorm}(x_l)
\right)
\]

\vspace{0.4cm}

Modern large models usually prefer Pre-Norm.

\end{frame}

%------------------------------------------------
\begin{frame}{RMSNorm}

RMSNorm simplifies LayerNorm by normalizing using the root mean square:

\[
\mathrm{RMS}(x)
=
\sqrt{
\frac{1}{d}
\sum_{i=1}^{d}x_i^2
}
\]

\[
\mathrm{RMSNorm}(x)
=
\gamma
\frac{x}{\mathrm{RMS}(x)}
\]

\vspace{0.4cm}

Used in many modern LLMs.

\end{frame}

%------------------------------------------------
\section{Attention Approximation}

\begin{frame}{The Cost of Full Attention}

Full self-attention has quadratic complexity:

\[
O(n^2)
\]

where $n$ is sequence length.

\vspace{0.4cm}

This becomes expensive for:

\begin{itemize}
    \item Long documents
    \item Long conversations
    \item Code repositories
    \item Multimodal inputs
\end{itemize}

\end{frame}

%------------------------------------------------
\begin{frame}{Sparse Attention}

Instead of allowing every token to attend to every other token:

\[
n \times n
\]

we restrict attention to selected patterns.

\vspace{0.4cm}

Examples:

\begin{itemize}
    \item Local attention
    \item Sliding window attention
    \item Global tokens
    \item Longformer-style sparse attention
\end{itemize}

\end{frame}

%------------------------------------------------
\begin{frame}{Sliding Window Attention}

In Sliding Window Attention, each token attends only to nearby tokens.

\vspace{0.4cm}

\[
\mathrm{Attention}(i)
=
\{j : |i-j| \leq w\}
\]

\vspace{0.4cm}

where $w$ is the window size.

\vspace{0.4cm}

This is similar to a receptive field in convolutional neural networks.

\end{frame}

%------------------------------------------------
\begin{frame}{Local and Global Attention}

Some models combine:

\begin{itemize}
    \item Local attention for nearby context
    \item Global attention for special tokens
\end{itemize}

\vspace{0.4cm}

This allows:

\begin{itemize}
    \item Better scalability
    \item Long-document processing
    \item Reduced memory cost
\end{itemize}

\end{frame}

%------------------------------------------------
\begin{frame}{Sharing Attention Heads}

Standard Multi-Head Attention uses separate $Q,K,V$ for each head.

\vspace{0.4cm}

Modern models often share key/value heads:

\begin{itemize}
    \item MHA: Multi-Head Attention
    \item MQA: Multi-Query Attention
    \item GQA: Grouped-Query Attention
\end{itemize}

\vspace{0.4cm}

This reduces inference memory and improves decoding speed.

\end{frame}

%------------------------------------------------
\begin{frame}{MHA, MQA, GQA}

\begin{table}
\centering
\begin{tabular}{lll}
\toprule
Method & Key/Value Sharing & Typical Use \\
\midrule
MHA & No sharing & Original Transformer \\
MQA & One KV pair for all heads & Fast inference \\
GQA & KV shared by groups & Modern LLMs \\
\bottomrule
\end{tabular}
\end{table}

\vspace{0.4cm}

GQA is a compromise between quality and efficiency.

\end{frame}

%------------------------------------------------
\section{Transformer-Based Models}

\begin{frame}{Three Families of Transformer Models}

\begin{table}
\centering
\begin{tabular}{lll}
\toprule
Family & Main Use & Examples \\
\midrule
Encoder-Decoder & Text-to-text & T5, mT5, ByT5 \\
Encoder-Only & Classification, understanding & BERT, RoBERTa \\
Decoder-Only & Text generation & GPT, Llama \\
\bottomrule
\end{tabular}
\end{table}

\end{frame}

%------------------------------------------------
\begin{frame}{Encoder-Decoder Models}

Encoder-decoder models map input text to output text.

\vspace{0.3cm}

Examples:

\begin{itemize}
    \item Machine translation
    \item Summarization
    \item Question answering
\end{itemize}

\vspace{0.4cm}

Representative models:

\begin{itemize}
    \item T5
    \item mT5
    \item ByT5
\end{itemize}

\end{frame}

%------------------------------------------------
\begin{frame}{Encoder-Only Models}

Encoder-only models produce contextual representations.

\vspace{0.4cm}

Typical tasks:

\begin{itemize}
    \item Sentiment classification
    \item Named entity recognition
    \item Semantic similarity
    \item Text classification
\end{itemize}

\vspace{0.4cm}

Examples:

\begin{itemize}
    \item BERT
    \item RoBERTa
    \item DistilBERT
\end{itemize}

\end{frame}

%------------------------------------------------
\begin{frame}{Decoder-Only Models}

Decoder-only models predict the next token.

\[
P(x_t \mid x_1,\ldots,x_{t-1})
\]

\vspace{0.4cm}

Examples:

\begin{itemize}
    \item GPT family
    \item Llama
    \item Mistral
    \item Qwen
\end{itemize}

\vspace{0.4cm}

This architecture dominates modern generative LLMs.

\end{frame}

%------------------------------------------------
\section{BERT}

\begin{frame}{Why BERT?}

\[
\text{BERT}
=
\text{Bidirectional Encoder Representations from Transformers}
\]

\vspace{0.4cm}

Key ideas:

\begin{itemize}
    \item Encoder-only Transformer
    \item Bidirectional self-attention
    \item Pretraining followed by fine-tuning
\end{itemize}

\end{frame}

%------------------------------------------------
\begin{frame}{BERT is Bidirectional}

Unlike decoder-only models, BERT can attend to both left and right context.

\vspace{0.4cm}

Example:

\begin{quote}
The student went to the bank to deposit money.
\end{quote}

\vspace{0.4cm}

The meaning of ``bank'' depends on surrounding context.

\end{frame}

%------------------------------------------------
\begin{frame}{BERT Strategy}

BERT uses a two-stage training process:

\begin{enumerate}
    \item Pretraining with proxy tasks
    \item Fine-tuning on a downstream task
\end{enumerate}

\vspace{0.4cm}

Pretraining tasks:

\begin{itemize}
    \item Masked Language Modeling (MLM)
    \item Next Sentence Prediction (NSP)
\end{itemize}

\end{frame}

%------------------------------------------------
\begin{frame}{Masked Language Modeling}

Input:

\begin{quote}
The capital of France is [MASK].
\end{quote}

\vspace{0.4cm}

Target:

\[
\text{Paris}
\]

\vspace{0.4cm}

The model learns to predict masked tokens using both left and right context.

\end{frame}

%------------------------------------------------
\begin{frame}{BERT Input Processing}

BERT input includes:

\begin{itemize}
    \item WordPiece tokens
    \item Special token \texttt{[CLS]}
    \item Separator token \texttt{[SEP]}
    \item Segment embeddings
    \item Position embeddings
\end{itemize}

\vspace{0.4cm}

Example:

\[
\texttt{[CLS] A cute teddy bear is reading . [SEP]}
\]

\end{frame}

%------------------------------------------------
\begin{frame}{Fine-Tuning BERT}

After pretraining, BERT can be fine-tuned for:

\begin{itemize}
    \item Classification
    \item Question answering
    \item Natural language inference
    \item Token classification
\end{itemize}

\vspace{0.4cm}

Often, only a small task-specific head is added.

\end{frame}

%------------------------------------------------
\begin{frame}{BERT: Pros and Cons}

\begin{table}
\centering
\begin{tabular}{ll}
\toprule
Pros & Cons \\
\midrule
Strong language understanding & Not ideal for generation \\
Needs less labeled data & Requires fine-tuning \\
Bidirectional context & Less flexible than modern LLMs \\
\bottomrule
\end{tabular}
\end{table}

\end{frame}

%------------------------------------------------
\begin{frame}{BERT Variants}

Important variants:

\begin{itemize}
    \item RoBERTa
    \item DistilBERT
    \item ALBERT
\end{itemize}

\vspace{0.4cm}

These variants improve training efficiency, compression, or performance.

\end{frame}

%------------------------------------------------
\begin{frame}{Key Takeaways}

\begin{itemize}
    \item Position information is essential for Transformers.
    \item RoPE is widely used in modern LLMs.
    \item LayerNorm and RMSNorm stabilize deep Transformer training.
    \item Sparse attention and KV sharing improve efficiency.
    \item Transformer models come in encoder-decoder, encoder-only, and decoder-only families.
    \item BERT is an encoder-only model focused on language understanding.
\end{itemize}

\end{frame}

%------------------------------------------------
\begin{frame}{Suggested Reading}

\begin{itemize}
    \item Vaswani et al. (2017), \textit{Attention Is All You Need}
    \item Devlin et al. (2018), \textit{BERT: Pre-training of Deep Bidirectional Transformers}
    \item Su et al. (2021), \textit{RoFormer: Enhanced Transformer with Rotary Position Embeddings}
    \item Press et al. (2021), \textit{Train Short, Test Long: Attention with Linear Biases}
    \item Xiong et al. (2020), \textit{On Layer Normalization in the Transformer Architecture}
    \item Beltagy et al. (2020), \textit{Longformer}
\end{itemize}

\end{frame}

%------------------------------------------------
\begin{frame}
\centering
\Huge Questions?
\end{frame}

\end{document}